\section{Validate a Binary Search Tree}
\label{sec:bst-validate}

\problemstatement{Given the root of a binary tree, determine whether it
is a valid BST: for every node, all values in its left subtree must be
strictly less than the node's value, and all values in its right subtree
must be strictly greater.}

\begin{patternbox}
The keyword trap here: it is tempting to check only that each node is
greater than its immediate left child and less than its immediate right
child, but the BST invariant is actually about \textbf{every} node in
each subtree, not just the direct children. \emph{``Validate a property
that must hold against \textbf{all ancestors}, not just the immediate
parent''} signals that the recursion needs to carry a \textbf{valid
range} downward, narrowing at every step, rather than just comparing
parent to child.
\end{patternbox}

\subsection*{Understanding the problem carefully}

Consider the tree with root $5$, left child $3$, and $3$'s right child
$6$. Checking only parent-child pairs would see: $3<5$ (fine, left child
smaller than parent) and $6>3$ (fine, right child larger than its own
parent) -- and incorrectly conclude this is a valid BST. But $6$ sits in
$5$'s \emph{left} subtree, and the BST invariant requires \emph{every}
value there to be less than $5$, which $6$ violates. The bug is checking
only local parent-child relationships instead of the full inherited
range of validity.

\begin{ideabox}[Pass Down a Valid (Lower, Upper) Range]
Recurse with two bounds, $\textit{lower}$ and $\textit{upper}$
(initially $-\infty$ and $+\infty$ at the root). At each node: if its
value is not strictly between $\textit{lower}$ and $\textit{upper}$,
the tree is invalid. Otherwise, recurse left with the same
$\textit{lower}$ but $\textit{upper}$ tightened to the current node's
value (everything in the left subtree must be less than this node), and
recurse right with $\textit{lower}$ tightened to the current node's
value and the same $\textit{upper}$.
\end{ideabox}

\subsection*{Why narrowing the range at every level correctly enforces the global invariant}

By construction, the range $(\textit{lower}, \textit{upper})$ passed into
any node already reflects the combined constraints of \emph{every}
ancestor on the path from the root, not just the immediate parent: each
recursive call only ever tightens the range, never loosens it, so a deep
descendant's range is the \emph{intersection} of every constraint imposed
by every ancestor above it. Checking the current node's value against
this fully-accumulated range is therefore equivalent to checking it
against every single ancestor directly, just done incrementally and in
$O(1)$ work per node instead of re-walking ancestors at every step.

\subsection*{Algorithm}

\begin{enumerate}
  \item \textsc{Validate}(\textit{node}, \textit{lower}, \textit{upper}):
  \begin{enumerate}
    \item If $\textit{node} = $ null: return true (an empty tree is
          trivially valid).
    \item If $\textit{node}.\textit{val} \le \textit{lower}$ or
          $\textit{node}.\textit{val} \ge \textit{upper}$: return
          false.
    \item Return
          \textsc{Validate}($\textit{node}.\textit{left}, \textit{lower},
          \textit{node}.\textit{val}$) \textbf{and}
          \textsc{Validate}($\textit{node}.\textit{right},
          \textit{node}.\textit{val}, \textit{upper}$).
  \end{enumerate}
  \item Call with $\textit{lower}=-\infty$, $\textit{upper}=+\infty$ at
        the root.
\end{enumerate}

\subsection*{Dry run}

\dryrun{Take the tree with root $5$, left child $3$, $3$'s right child
$6$.}

\begin{itemize}
  \item At $5$, range $(-\infty,+\infty)$: $5$ is within range.
        Recurse left into $3$ with range $(-\infty,5)$; recurse right
        into null with range $(5,+\infty)$ (trivially valid).
  \item At $3$, range $(-\infty,5)$: $3$ is within range. Recurse left
        into null (valid); recurse right into $6$ with range $(3,5)$.
  \item At $6$, range $(3,5)$: is $6$ strictly between $3$ and $5$? No
        ($6 \ge 5$). Return false.
\end{itemize}

The tree is correctly identified as \textbf{invalid}.

\subsection*{C++ implementation}

\begin{lstlisting}
#include <bits/stdc++.h>
using namespace std;

struct TreeNode {
    int val;
    TreeNode* left;
    TreeNode* right;
    TreeNode(int v) : val(v), left(nullptr), right(nullptr) {}
};

bool validate(TreeNode* node, long long lower, long long upper) {
    if (node == nullptr) return true;
    if (node->val <= lower || node->val >= upper) return false;

    return validate(node->left, lower, node->val) &&
           validate(node->right, node->val, upper);
}

bool isValidBST(TreeNode* root) {
    return validate(root, LLONG_MIN, LLONG_MAX);
}

int main() {
    TreeNode* root = new TreeNode(5);
    root->left = new TreeNode(3);
    root->left->right = new TreeNode(6);

    cout << (isValidBST(root) ? "true" : "false") << endl;
    return 0;
}
\end{lstlisting}

\begin{complexitybox}
\textbf{Time Complexity.} Every node is visited exactly once, with
$O(1)$ work per node, giving
\[
  O(n).
\]
\textbf{Space Complexity.} $O(h)$ for the recursion stack.
\end{complexitybox}

\begin{takeawaybox}
Any property phrased as ``true for every node relative to \emph{all} of
its ancestors/descendants,'' not just its immediate neighbors, needs
state (here, a valid range) threaded through the recursion, accumulating
constraints as it descends. An equally valid alternative solution
performs an inorder traversal and checks that the resulting sequence is
strictly increasing, directly invoking
Section~\ref{sec:bst-intro}'s headline fact -- worth knowing as a second
way to arrive at the same answer.
\end{takeawaybox}
